% Rigidity of the rank-48 4x4 matrix-multiplication scheme under
% solved flip moves.  Build: tectonic doc/matmul_rigid48_paper.tex
\documentclass[10pt]{article}
\usepackage[margin=2.4cm]{geometry}
\usepackage{amsmath,amssymb}
\usepackage{booktabs}
\usepackage{array}
\usepackage{fancyvrb}
\usepackage{enumitem}
\usepackage{microtype}
\usepackage[hidelinks]{hyperref}
\urlstyle{same}
\setlist{itemsep=2pt,topsep=3pt}
\setlength{\emergencystretch}{3em}
\DefineVerbatimEnvironment{Code}{Verbatim}{fontsize=\footnotesize}

\title{\bfseries Rigidity of the Rank-48 $4\times4$ Matrix-Multiplication
  Scheme under Solved Flip Moves --- a Machine-Checked Obstruction}
\author{Greg Sidebottom}
\date{}

\begin{document}
\maketitle
\vspace{-3em}
\begin{center}
\itshape Research note --- logic repo
(\url{https://github.com/gsidebottom/logic}), matmul track, 2026-07-08.
Companion to ``A 55-Addition Rank-23 Scheme for $3\times3$ Matrix
Multiplication via Exact Two-Sided Minimization''
(\texttt{doc/matmul\_adds\_paper.md}, artifacts
DOI~\path{10.5281/zenodo.21240904}); every claim here has a mechanical
check (\S\ref{sec:repro}), and the main theorem is machine-checked in
Lean~4 (\S\ref{sec:lean}).
\end{center}

\section*{Abstract}

Fast matrix multiplication of $4\times4$ matrices currently stands at
\textbf{48 multiplications} over the real numbers --- a scheme found by
AlphaEvolve (2025) over the complex numbers and given rational
coefficients by Dumas, Pernet, and Sedoglavic (2025). Whether 47 is
possible outside characteristic~2 is a central open problem, and the
most productive discovery machinery of recent years --- \textbf{flip
graphs} (Kauers--Moosbauer 2022) --- has never produced a
characteristic-0 improvement at this format despite sustained effort.

This note offers a certified explanation. We generalize flip moves
from their native mod-2 setting to the rationals, where the flip
scalar becomes a free parameter $\lambda$ and two new phenomena
appear: exact factor coincidences --- the events that enable rank
reductions --- occur with probability zero under random moves, so
productive moves must be \textbf{solved for} (a two-unknown linear
system per candidate); and solvable ``coincidence'' moves turn out to
be abundant. We then prove, by exhaustive enumeration in exact
arithmetic:

\begin{quote}
\textbf{Theorem (rigidity).} From the rational rank-48 scheme, the
graph generated by one \emph{split} (any pair of summands, any tensor
slot) followed by arbitrary sequences of \emph{solved flips} closes on
a finite component of exactly \textbf{7{,}408 states}, every state of
which is a verified decomposition of the $4\times4$ multiplication
tensor, and \textbf{no state admits a reduction below 49 summands
except the trivial return to the seed}. In particular, no rank-47
scheme --- which would be a characteristic-0 record --- is reachable
this way.
\end{quote}

The theorem is verified end to end in \textbf{Lean~4} by reflection
(\texttt{native\_decide}), with deduplication on full canonical forms
(no hashes), unbounded exact integers (no magnitude caps), saturation
of the search asserted \emph{inside} the decided proposition, and the
seed self-certified against the definition of the multiplication
tensor.

\section{Background, with definitions}\label{sec:background}

\paragraph{Bilinear schemes and rank.}
Multiplying two $n\times n$ matrices is a bilinear map. A
\emph{bilinear scheme} (equivalently, a rank-$r$ decomposition of the
multiplication tensor $\langle n,n,n\rangle$) computes it as
\[
  M_t=\Bigl(\sum_{i,j}\alpha_t[i,j]\,A[i,j]\Bigr)\cdot
      \Bigl(\sum_{k,l}\beta_t[k,l]\,B[k,l]\Bigr),\quad t=1..r,\qquad
  C[p,q]=\sum_t \gamma_t[p,q]\,M_t
\]
for coefficient vectors $(\alpha_t,\beta_t,\gamma_t)$. Each product
$M_t$ is a \emph{rank-one summand} $a_t\otimes b_t\otimes c_t$ of the
tensor; $r$ is the scheme's \textbf{rank} --- the number of
(expensive, non-commutative) multiplications. Validity is a system of
polynomial constraints on the coefficients, the \textbf{Brent
equations} (Brent 1970): for $\langle 4,4,4\rangle$ there are
$16^3=4096$ of them. Because the products never commute $A$- and
$B$-entries, a valid scheme works over any ring --- in particular over
\emph{block matrices}, which is what makes low rank matter: a rank-$r$
scheme for $\langle4,4,4\rangle$ yields an $O(n^{\log_4 r})$ algorithm
by recursion.

\paragraph{Characteristic.}
The \emph{characteristic} of a field is the smallest number of times
$1$ must be added to itself to reach $0$ --- e.g.\ $2$ for the Boolean
field $\mathrm{GF}(2)=\{0,1\}$ (where $1+1=0$), and
\emph{characteristic 0} for $\mathbb{Q},\mathbb{R},\mathbb{C}$, where
no repeated sum of $1$'s vanishes. The distinction is essential here:
a scheme with coefficients in $\mathrm{GF}(2)$ computes matrix
products \emph{only} for matrices over fields of characteristic~2.
AlphaTensor's celebrated rank-47 scheme for $\langle4,4,4\rangle$
(Fawzi et al.\ 2022) is of this kind --- it does \textbf{not} give a
47-multiplication algorithm for real matrices. For real (or
floating-point, or integer) linear algebra one needs
characteristic-0 coefficients, and there the record is:

\begin{center}
\begin{tabular}{@{}lll@{}}
\toprule
rank & scheme & coefficients\\
\midrule
49 & Strassen 1969, squared ($2\times2$ scheme applied twice) & $\pm1$\\
48 & AlphaEvolve (Novikov et al.\ 2025) & $\mathbb{C}$ (halves of Gaussian integers)\\
48 & Dumas--Pernet--Sedoglavic (DPS) rational point of the same scheme & $\mathbb{Z}[\tfrac12]$\\
47 & \textbf{open} in characteristic 0 & ---\\
\bottomrule
\end{tabular}
\end{center}

Two structural facts sharpen the picture (Moran--Schwartz--Yuan 2026):
the AlphaEvolve scheme admits \textbf{no integer-coefficient
equivalent} (the halves are intrinsic), and the only other known
rank-48 scheme --- Kaporin's (2024), found independently by nonlinear
least squares --- admits \textbf{no real-coefficient equivalent at
all} (its orbit requires cube roots of 2). So the DPS rational scheme
is, up to the natural symmetries, \emph{the} rank-48 object available
to characteristic-0 computation, and the question ``can flip methods
improve $4\times4$ over the reals?'' is a question about this one
scheme. We call it \textbf{the seed} throughout.

\paragraph{Equivalence (de Groote symmetry).}
Two schemes are \emph{equivalent} if one maps to the other under the
isotropy group of the multiplication tensor --- invertible
``sandwiching'' of the three coefficient families plus permutations of
the three tensor slots (de Groote 1978). All statements about ``the''
scheme are statements about its orbit.

\paragraph{Flip graphs.}
Kauers and Moosbauer (2022) organize all rank-$r$ schemes into a
graph. Two summands that share a factor in one slot, say
$a_i=a_j=a$, admit a \textbf{flip}:
\[
  a\otimes b_i\otimes c_i + a\otimes b_j\otimes c_j
  \;=\;
  a\otimes(b_i{+}b_j)\otimes c_i + a\otimes b_j\otimes(c_j{-}c_i)
\]
--- an exchange of material that preserves the sum (hence validity)
and the rank. Two summands proportional in \emph{two} slots merge
into one --- a \textbf{reduction}, decreasing the rank. A
\textbf{split} (Moosbauer--Poole 2025) is the reverse of a reduction:
replace a summand by two whose chosen-slot factors sum to it,
\emph{increasing} the rank to escape local structure. Random walks
over flips punctuated by reductions found new records in
characteristic 2 (among them $5\times5$ in 93 multiplications),
making flip graphs the most productive scheme-discovery machinery
outside large-scale machine learning.

\paragraph{The failures this note explains.}
At the $4\times4$ format the flip approach has conspicuously stalled,
in ways we and others have measured:
\begin{itemize}
\item In characteristic 2, the two known rank-47 schemes
  (AlphaTensor's and the flip-graph one) sit in near-isolated
  flip-graph components; our earlier experiments (engine
  \texttt{src/flip.rs} in the companion repository) additionally
  verified that \textbf{neither lifts to $\pm1$ integer coefficients}
  --- so neither yields a characteristic-0 algorithm --- and that
  mod-2 flip descent from random rank-49 starting points stalls: our
  probe of ${\sim}100{,}000$ descent runs found every rank-49 landing
  point to be \textbf{100\% flip-isolated} (no two summands sharing
  any factor --- a flip-graph sink).
\item In characteristic 0 the machinery could not even be applied to
  the one interesting seed: the DPS scheme's coefficients contain
  $\tfrac12$'s, so it \textbf{has no mod-2 reduction at all}, and the
  published flip tooling is mod-$p$.
\end{itemize}
This note builds the missing rational tooling, measures what the
rational flip landscape around the seed actually looks like, and
proves that the natural search cannot leave it.

\section{Rational flip moves}\label{sec:moves}

Everything below is implemented in \texttt{src/bin/flip48.rs} (Rust),
with exact arithmetic throughout: a factor vector is stored as sixteen
integers times a power of two (the coefficient ring
$\mathbb{Z}[\tfrac12]$ of the seed is closed under all moves), and
every walk state is verifiable against the full Brent system in
128-bit exact integer arithmetic.

\paragraph{Canonical gauge.}
Each summand $a\otimes b\otimes c$ carries a projective freedom
$(a,b,c)\to(ua,\,vb,\,(uv)^{-1}c)$. We fix it by keeping $a$ and $b$
\emph{primitive} --- integer entries with no common factor, leading
entry positive --- with all scalar content folded into $c$. Under
this gauge, proportionality of $a$- or $b$-factors is literal
equality, which is what the move machinery tests. (An early version
of the engine canonicalized only powers of two, making
proportionality by odd factors invisible --- a completeness bug worth
recording, since walks create factors of 3 immediately.)

\paragraph{$\lambda$-flips.}
Over a field the flip identity holds with any scalar:
\[
  a\otimes b_i\otimes c_i + a\otimes b_j\otimes c_j
  =
  a\otimes(b_i{+}\lambda b_j)\otimes c_i
  + a\otimes b_j\otimes(c_j{-}\lambda c_i).
\]
Over $\mathrm{GF}(2)$ the only choice is $\lambda=1$; over
$\mathbb{Q}$ the move is a one-parameter family. The engine restricts
$\lambda$ to $\pm2^k$ (dyadic scalars) so that all coefficients remain
in $\mathbb{Z}[\tfrac12]$.

\paragraph{Reductions on any two slots.}
Summands proportional in slots $(a,b)$, $(a,c)$, or $(b,c)$ merge on
the third slot, with the proportionality ratios folded in exactly.

\paragraph{Universal splits.}
Over $\mathbb{Q}$, splits acquire a power they lack mod 2:
\emph{any} summand $i$ can be split against \emph{any} other summand
$k$ in any slot --- write $a_i=\mu a_k+(a_i-\mu a_k)$ ---
manufacturing a shared factor with a partner of our choosing.
Mobility can be engineered; it need not be found.

\section{What the rational flip landscape looks like}\label{sec:landscape}

Three measurements, all exact, set up the theorem.

\textbf{(1) The seed is flip-isolated.} Its 48 summands are pairwise
non-proportional in \emph{every} slot: not a single classical flip is
available. This is the same signature as the characteristic-2
rank-47 schemes --- the interesting objects at this format all sit at
flip-graph sinks. Splits are therefore not an optimization but the
only way to move at all.

\textbf{(2) Over $\mathbb{Q}$, coincidences never happen by chance.}
A reduction requires two summands to become \emph{exactly}
proportional in two slots. Mod 2 the state space is finite and random
walks hit such coincidences constantly; over $\mathbb{Q}$ the walk
wanders an infinite space in which exact coincidence is a
measure-zero event. Empirically: ten million random-$\lambda$ moves
from the seed produced \textbf{zero} coincidences. Every reduction
ever observed was a split un-doing itself. Random walking --- the
engine of all published flip-graph results --- is the wrong paradigm
for characteristic 0.

\textbf{(3) Solvable coincidences are abundant.} The cure is to
\emph{solve} for the moves. For a pair $(i,j)$ sharing slot $s$ and a
third summand $m$, the requirement ``after a flip with scalar
$\lambda$, factor $f_i$ becomes proportional to $f_m$'' is the
16-equation, 2-unknown linear system
\[
  f_i+\lambda f_j=\mu f_m,
\]
solvable exactly by Cramer's rule; the engine accepts solutions with
dyadic $\lambda$. Such \emph{coplanar triples} turn out to be
plentiful: around 3.5 million solved flips were applied in a
45-second, 12-thread run. The search space is rich --- which makes
what follows meaningful.

\section{The rigidity theorem}\label{sec:theorem}

With mobility engineered by splits and productive moves solved for,
the reachable landscape can be enumerated outright --- no randomness.

\paragraph{Enumeration.}
For every ordered pair of summands $(i,k)$ and every slot $s$ ---
$48\cdot47\cdot3=6768$ roots --- split $i$ against $k$ in $s$; then
explore by depth-first search: at each state, find every
shared-factor pair, compute every solved flip, apply each, and
deduplicate states; at every state, run the reduction sweep and
record the outcome. The split scalar $\mu$ does not appear in any
coincidence system (the split part inherits the other two factors of
the split summand unchanged), so taking $\mu=1$ loses no generality
for reachability of coincidences.

\paragraph{Result.}
The exploration \textbf{saturates}: the union of all components
closes after ${\sim}8500$ node visits (7{,}408 distinct states under
global deduplication), stable from depth 4 onward and identical under
magnitude caps 256 and 1024 (the component never approaches either).
Within it:
\begin{itemize}
\item every state is a valid rank-49 or rank-48 decomposition
  (verified against all 4096 Brent equations, exactly);
\item \textbf{every reduction sweep terminates at $\ge49$ summands,
  or at 48 by the trivial merge that returns the original seed};
\item in particular, \textbf{no rank-47 state exists in the
  component}, and no rank-48 state other than the seed --- the
  solved-move system cannot even produce a \emph{different}
  presentation of rank 48, let alone a smaller one.
\end{itemize}
We call this property \textbf{rigidity of the seed under solved
moves}. The scope is stated precisely: one split (all pairs, all
slots, all dyadic scalars via the $\mu$-independence above), then
arbitrary sequences of solved flips with dyadic $\lambda$; shared
factors in the two canonical slots are detected up to
proportionality, in the scalar-carrying slot by exact equality. What
lies outside the scope --- and outside any current method's reach ---
is discussed in \S\ref{sec:discussion}.

\section{The machine-checked proof}\label{sec:lean}

Exhaustive computational claims deserve mechanical verification; the
enumeration above is small enough to re-run inside a proof assistant.
The Lean~4 development
\texttt{proof/Mm/Rigid48.lean} (same Lake project
as the $3\times3$ correctness proof of the companion paper)
re-implements the entire move system --- gauge, $\lambda$-flips,
splits, all three reduction patterns, the Cramer solver, the DFS ---
as \emph{total} executable functions over Lean's unbounded integers,
and proves:
\begin{Code}
theorem seed_valid : checkDecomp seed = true
theorem rigid : explore seed 20000 = <7408, true, true, true>
\end{Code}
by reflection (\texttt{native\_decide}): the enumeration runs inside
Lean's native evaluator and the kernel accepts the resulting boolean.
The four components of the pinned tuple assert, in order: the
component has exactly 7{,}408 states; the search \emph{saturated}
(fuel was not exhausted --- closure is part of the theorem, not an
external observation); every visited state passes the full Brent
check; and no reduction sweep lands below 49 summands except at the
seed itself.

The formalization is stronger than the Rust run in four ways:
\begin{enumerate}
\item \textbf{collision-free deduplication} --- visited states are
  keyed by their full canonical forms, not 64-bit hashes;
\item \textbf{no magnitude caps} --- Lean's arbitrary-precision
  integers remove the caps entirely from the statement;
\item \textbf{saturation inside the decided proposition};
\item \textbf{a self-certifying seed} --- the 48-summand literal
  (exported by \texttt{flip48 --emit-lean}) is proved to decompose
  the \emph{defined} multiplication tensor, so no trust in the
  exporter is required.
\end{enumerate}

Trust base: Lean's standard axioms plus the \texttt{native\_decide}
compiler-trust axiom (printed by the build's \texttt{\#print axioms}
audit). The component count was first learned from a compiled probe
executable and only then pinned into the theorem --- a wrong pin
fails the build, guarding against a vacuous pass.
\texttt{lake exe rigid48probe} re-runs the exploration natively in
about $3\tfrac12$ minutes.

\section{What rigidity does and does not mean}\label{sec:discussion}

\paragraph{It explains the observed failures.}
The one rank-48 scheme available to characteristic-0 computation sits
at a flip sink; random rational walks provably cannot generate
coincidences; and the systematic solved-move search closes on a small
component containing no exit. Every mechanism by which flip methods
have ever found records is individually and jointly obstructed at
this seed --- consistent with three years of silence at $4\times4$ in
characteristic 0, and with our own measured stalls in
characteristic~2.

\paragraph{It is scoped, honestly.}
Rigidity is proved for the solved-move system: single splits, dyadic
$\lambda$, solved flips. Conceivable escapes outside the scope:
multi-split states (rank 50+ before descending); non-dyadic rational
$\lambda$ (leaving $\mathbb{Z}[\tfrac12]$; the seed's own ring makes
this unnatural but not impossible); moves through \emph{unsolved}
territory followed by later solves (the measure-zero argument makes
this unpromising but it is heuristic, not proven); and entirely
different seeds --- though the classification results quoted in
\S\ref{sec:background} say there are none over $\mathbb{R}$ at rank
48 today. Extending the certificate to two-split radii is a finite
(larger) enumeration of the same kind.

\paragraph{A structural reading.}
The Moran--Schwartz--Yuan results show the scheme's
$\tfrac12$-coefficients are intrinsic; our measurements show its
factor geometry is maximally spread (48 distinct directions in every
slot, coplanarities present but never twice-aligned). Rigidity adds:
the scheme is not merely isolated but \emph{locally unique} --- its
neighborhood in the solved-move graph contains no second rank-48
point. Discovery of a characteristic-0 rank-47, if one exists, will
require machinery that does not move through this landscape ---
large-scale learned search (how the 48 itself was found), algebraic
construction, or SAT-style exhaustive methods at presently
inaccessible scales.

\section{Reproduction}\label{sec:repro}

Prerequisites: a Rust toolchain; \texttt{elan} for the Lean part
(toolchain and Mathlib commit pinned in the repository).

\begin{Code}
# the engine: exact seed verification + the walk/enumeration findings
cargo build --release --bin flip48
./target/release/flip48 --seconds 45 --threads 12      # walk stats
./target/release/flip48 --pursue3 10 --cap 256         # saturation
./target/release/flip48 --pursue3 10 --cap 1024        # cap-independence

# the machine-checked theorem
cd proof
lake exe cache get     # prebuilt Mathlib
lake build             # elaborates rigid + seed_valid (native_decide;
                       # the Rigid48 module takes ~15 minutes)
lake exe rigid48probe  # re-runs the exploration, prints
                       # visited=7408 saturated=true allValid=true noBad=true
\end{Code}

The seed's provenance: \texttt{matmul/dps48/\{L,R,P\}.sms} are the
Dumas--Pernet--Sedoglavic matrices fetched from their PLinOpt
repository; \texttt{matmul/probe48.py verify} proves them against all
4096 Brent equations over exact rationals; \texttt{flip48
--emit-lean} exports the gauged seed, which Lean re-certifies
independently (\texttt{seed\_valid}).

\section*{Acknowledgments}

This work was carried out in an extended interactive collaboration
with Claude (Anthropic; the Fable~5 and Opus~4.8 models), which
implemented the search engines, the Lean formalization, and much of
this text under the author's direction and review. Every
computational claim is mechanically checkable by the commands in
\S\ref{sec:repro}; the results rest on those independent
verifications rather than on trust in either the author or the tools.

\section*{References}

\begin{itemize}
\item V. Strassen. \emph{Gaussian Elimination is not Optimal.}
  Numerische Mathematik 13:354--356 (1969).
\item R. P. Brent. \emph{Algorithms for Matrix Multiplication.}
  Stanford report CS-TR-70-157 (1970).
\item H. F. de Groote. \emph{On Varieties of Optimal Algorithms for
  the Computation of Bilinear Mappings I/II.} Theoretical Computer
  Science 7 (1978).
\item A. Fawzi, M. Balog, A. Huang, T. Hubert, B. Romera-Paredes,
  M. Barekatain, A. Novikov, F. J. R. Ruiz, J. Schrittwieser,
  G. Swirszcz, D. Silver, D. Hassabis, P. Kohli. \emph{Discovering
  Faster Matrix Multiplication Algorithms with Reinforcement
  Learning} (AlphaTensor). Nature 610:47--53 (2022).
\item M. Kauers, J. Moosbauer. \emph{Flip Graphs for Matrix
  Multiplication.} arXiv:2212.01175 (2022).
\item J. Moosbauer, M. Poole. \emph{Flip-Graph Algorithms for Matrix
  Multiplication in Larger Formats} (and successors); see also
  M. Kauers, I. Wood, \emph{Exploring the Meta Flip Graph for Matrix
  Multiplication}, arXiv:2510.19787 (2025).
\item A. Novikov, N. V\~u, M. Eisenberger, E. Dupont, P.-S. Huang,
  A. Z. Wagner, S. Shirobokov, B. Kozlovskii, F. J. R. Ruiz,
  A. Mehrabian, M. P. Kumar, A. See, S. Chaudhuri, G. Holland,
  A. Davies, S. Nowozin, P. Kohli, M. Balog. \emph{AlphaEvolve: a
  Coding Agent for Scientific and Algorithmic Discovery.}
  arXiv:2506.13131 (2025).
\item J.-G. Dumas, C. Pernet, A. Sedoglavic. \emph{A Non-Commutative
  Algorithm for Multiplying $4\times4$ Matrices Using 48 Non-Complex
  Multiplications.} arXiv:2506.13242 (2025). Artifacts in the
  PLinOpt repository (\path{github.com/jgdumas/plinopt}).
\item J.-G. Dumas, C. Pernet, A. Sedoglavic. \emph{A More Accurate
  Rational Non-Commutative Algorithm for Multiplying $4\times4$
  Matrices Using 48 Multiplications.} arXiv:2603.18699 (2026).
\item I. Kaporin. \emph{Finding Complex-Valued Solutions of Brent
  Equations Using Nonlinear Least Squares.} Computational Mathematics
  and Mathematical Physics 64(9):1881--1891 (2024).
\item Y. Moran, O. Schwartz, S. Yuan. \emph{Complex to Rational Fast
  Matrix Multiplication.} arXiv:2602.13171 (2026).
\item G. Sidebottom. \emph{A 55-Addition Rank-23 Scheme for
  $3\times3$ Matrix Multiplication via Exact Two-Sided Minimization.}
  Companion paper, this repository (2026). Artifacts:
  DOI~\path{10.5281/zenodo.21240904}.
\item L. de Moura, S. Ullrich. \emph{The Lean 4 Theorem Prover and
  Programming Language.} CADE-28 (2021); and The Mathlib Community,
  \emph{The Lean Mathematical Library}, CPP (2020).
\end{itemize}

\end{document}
